% Строчки для суммарной таблицы, в следующей строке прописаны используемые экземпляры узлов
%===s control, LLN0, LVCBPTRC1, LVCBRECRBRE1, LVBTSRBLC1

%\needspace{4\baselineskip}
\color{uniblue}{\subsection{Логика отключения выключателя стороны НН трансформатора (ЛО НН)}}
\color{black}

\begin{enumerate}[label=\arabic{section}.\arabic{subsection}.\arabic{enumi}, labelsep=4pt, leftmargin=0pt, itemindent=57pt, itemsep=0pt, parsep=5pt] % если что можно крутить itemsep

%%%%%%%%%%%%%%%%%%%%%%%%%%%%%%%%%%%%%%%%%%%%%%%%%%%%%%%%%%%%%%%%%%%%% ОБЩИЕ СВЕДЕНИЯ %%%%%%%%%%%%%%%%%%%%%%%%%%%%%%%%%%%%%%%%%%%%%%%%%%%%%%%%%%%%%%

%\color{uniblue}{\subsubsection{Общие сведения}}
%\color{black}

\item Общие сведения

\begin{enumerate}[label=\arabic{section}.\arabic{subsection}.\arabic{enumi}.\arabic*, labelsep=4pt, leftmargin=0em, itemindent=65pt, parsep=0pt]

% ЗДЕСЬ ВЕТВЛЕНИЕ, если исполнение Т2 то вставляем первый абзац
\ifboolexpr{bool {isT2}}
{
    \item
    В составе устройства предусматриваются две функции <<ЛО НН1>> и <<ЛО НН2>> (для стороны НН1 и НН2 соответственно), ниже приводится общее описание работы функции <<ЛО НН>>. 
}
{}

\item
Функциональный блок <<ЛО НН>> служит для формирования команды отключения в схему выключателя ввода стороны НН трансформатора (функция <<ЛО>>), команды запрета АПВ (функция <<ЗАПВ>>) в схему АПВ выключателя ввода секции НН и команды запрета АВР (функция <<ЗАВР>>) в схему АВР секций НН. 
\item
В составе устройства предусматривается одна логика отключения стороны НН трансформатора (для схем управления вводным выключателелем секции НН).
\item
Функциональный блок <<ЛО НН>> может быть выведен из работы оперативно путем активации сигнала <<ОВ~ЛО~НН>>.

\end{enumerate}

%%%%%%%%%%%%%%%%%%%%%%%%%%%%%%%%%%%%%%%%%%%%%%%%%%%%%%%%%%%%%%%% ЛОГИКА ОТКЛЮЧЕНИЯ ВЫКЛЮЧАТЕЛЯ НН %%%%%%%%%%%%%%%%%%%%%%%%%%%%%%%%%%%%%%%%%%%%%%%%%%

%\color{uniblue}{\subsubsection{Логика отключения выключателя ввода стороны НН (ЛО)}}
%\color{black}

\item Логика отключения выключателя ввода стороны НН (ЛО)

\begin{enumerate}[label=\arabic{section}.\arabic{subsection}.\arabic{enumi}.\arabic*, labelsep=4pt, leftmargin=0em, itemindent=65pt, parsep=0pt]

\item
Ввод функции <<ЛО>> в работу на этапе параметрирования устройства осуществляется программным переключателем SGF1 <<Ввод\_функции>>. Информация о введенной функции отображается выходным сигналом <<ЛО НН / ЛО: Ввод>>.
\item
Функция <<ЛО>> может быть выведена из работы оперативно путем активации сигнала <<ОВ ЛО НН / ЛО>>, что характеризуется активным состоянием сигнала <<ЛО НН / ЛО: Оперативный вывод>> на выходе функции.  
\item
Функция принимает сигнал срабатывания от общей логики отключения трансформатора и формирует сигналы <<ЛО НН / ЛО: Отключение>> и <<ЛО НН / ЛО: Отключить аварийно>>, сигнал <<ЛО НН / ЛО: Отключить аварийно>> является импульсным, время импульса определяется выдержкой времени таймера T1 (уставка <<Тимп>>).
\item
Алгоритм работы <<ЛО>> выполнен в соответствии с рисунком \ref{lonn:img_lov}. В таблице \ref{lonn:tbl_lov} приведены параметры, необходимые для настройки функции <<ЛО>>.

\vspace{3mm}
\begin{figure}[H]
\centering
\includegraphics[width=0.9\textwidth,height=0.9\textheight,keepaspectratio]{\fbpath/61. ЛО В НН Т 35 Т2/_latex/img/lov.pdf}
\captionof{figure}{Функциональная схема <<ЛО>>}\label{lonn:img_lov}
\end{figure}

\small
\begin{longtable}{|>{\centering\arraybackslash}m{5cm}|>{\centering\arraybackslash}m{4cm}|>{\centering\arraybackslash}m{4cm}|>{\centering\arraybackslash}m{2cm}|>{\centering\arraybackslash}m{1cm}|}
\caption{Параметры для настройки функции <<ЛО>>\hfill\vspace{-0.5\baselineskip}}\label{lonn:tbl_lov}\\ 
\hline
\rowcolor{gray!30}
Параметр (Параметр на ИЧМ) & Условное обозначение на схеме & Значение/ Диапазон & Единица измерения & Шаг \\ 
\hline
\endfirsthead
%\caption{\emph{Продолжение\hfill\vspace{-0.5\baselineskip}}} \\ % Отступ обозначения таблицы от таблицы и выравнивание надписи слева
\caption*{\hspace{3pt}\emph{Продолжение таблицы \ref{lonn:tbl_lov}\hfill\vspace{-0.5\baselineskip}}} \\ % сделано по ГОСТ 2.105 п.6.8.7
\hline
\rowcolor{gray!30}
Параметр (Параметр на ИЧМ) & Условное обозначение на схеме & Значение/ Диапазон & Единица измерения & Шаг \\ 
%\hline
\endhead
%\hline
\endfoot
%\hline
\endlastfoot
%==+t1*LVCBPTRC1|LVTRESCBOFF1
%===t1*LVCBPTRC1
\centering Ввод функции в работу (Ввод\_функции) & \centering SGF1 & \centering 0 = Не предусмотрено\\1 = Предусмотрено & \centering -- & \centering \arraybackslash -- \\
\hline
\centering Длительность импульса (Тимп) & \centering T1 & \centering 0,05 ... 60,00 & \centering с & \centering \arraybackslash 0,01 \\
\hline
%===t1
\end{longtable}
\normalsize

\end{enumerate}
%%%%%%%%%%%%%%%%%%%%%%%%%%%%%%%%%%%%%%%%%%%%%%%%%%%%%%%%%%%%%%%% ЛОГИКА ЗАПРЕТА АПВ НН %%%%%%%%%%%%%%%%%%%%%%%%%%%%%%%%%%%%%%%%%%%%%%%%%%

%\color{uniblue}{\subsubsection{Логика запрета АПВ выключателя ввода стороны НН (Запрет АПВ)}}
%\color{black}

\item Логика запрета АПВ выключателя ввода стороны НН (Запрет АПВ)

\begin{enumerate}[label=\arabic{section}.\arabic{subsection}.\arabic{enumi}.\arabic*, labelsep=4pt, leftmargin=0em, itemindent=65pt, parsep=0pt]

\item
Ввод функции <<Запрет АПВ>> в работу на этапе параметрирования устройства осуществляется программным переключателем SGF1 <<Ввод\_функции>>. Информация о введенной функции отображается выходным сигналом <<ЛО НН / ЗАПВ: Ввод>>.
\item
Функция <<Запрет АПВ>> может быть выведена из работы оперативно путем активации сигнала <<ОВ ЛО НН / ЗАПВ>>, что характеризуется активным состоянием сигнала <<ЛО НН / ЗАПВ: Оперативный вывод>> на выходе функции.
\item
Функция принимает сигнал срабатывания от общей логики отключения и формирует сигналы <<ЛО НН / ЗАПВ: Запрет АПВ>>.
\item
Алгоритм работы <<Запрет АПВ>> выполнен в соответствии с рисунком \ref{lonn:img_zapv}. В таблице \ref{lonn:tbl_zapv} приведены параметры, необходимые для настройки функции <<Запрет АПВ>>.

\vspace{3mm}
\begin{figure}[H]
\centering
\includegraphics[width=0.9\textwidth,height=0.9\textheight,keepaspectratio]{\fbpath/61. ЛО В НН Т 35 Т2/_latex/img/zapv.pdf}
\captionof{figure}{Функциональная схема <<Запрет АПВ>>}\label{lonn:img_zapv}
\end{figure}

\small
\begin{longtable}{|>{\centering\arraybackslash}m{5cm}|>{\centering\arraybackslash}m{4cm}|>{\centering\arraybackslash}m{4cm}|>{\centering\arraybackslash}m{2cm}|>{\centering\arraybackslash}m{1cm}|}
\caption{Параметры для настройки функции <<Запрет АПВ>>\hfill\vspace{-0.5\baselineskip}}\label{lonn:tbl_zapv}\\ 
\hline
\rowcolor{gray!30}
Параметр (Параметр на ИЧМ) & Условное обозначение на схеме & Значение/ Диапазон & Единица измерения & Шаг \\ 
\hline
\endfirsthead
%\caption{\emph{Продолжение\hfill\vspace{-0.5\baselineskip}}} \\ % Отступ обозначения таблицы от таблицы и выравнивание надписи слева
\caption*{\hspace{3pt}\emph{Продолжение таблицы \ref{lonn:tbl_zapv}\hfill\vspace{-0.5\baselineskip}}} \\ % сделано по ГОСТ 2.105 п.6.8.7
\hline
\rowcolor{gray!30}
Параметр (Параметр на ИЧМ) & Условное обозначение на схеме & Значение/ Диапазон & Единица измерения & Шаг \\ 
%\hline
\endhead
%\hline
\endfoot
%\hline
\endlastfoot
%==+t1*LVCBRECRBRE1|LVTRESCBOFF1
%===t1*LVCBRECRBRE1
\centering Ввод функции в работу (Ввод\_функции) & \centering SGF1 & \centering 0 = Не предусмотрено\\1 = Предусмотрено & \centering -- & \centering \arraybackslash -- \\
\hline
%===t1
\end{longtable}
\normalsize

\end{enumerate}
%%%%%%%%%%%%%%%%%%%%%%%%%%%%%%%%%%%%%%%%%%%%%%%%%%%%%%%%%%%%%%%% ЛОГИКА ЗАПРЕТА АВР НН %%%%%%%%%%%%%%%%%%%%%%%%%%%%%%%%%%%%%%%%%%%%%%%%%%

%\color{uniblue}{\subsubsection{Логика запрета АВР стороны НН (Запрет АВР)}}
%\color{black}

\item Логика запрета АВР стороны НН (Запрет АВР)

\begin{enumerate}[label=\arabic{section}.\arabic{subsection}.\arabic{enumi}.\arabic*, labelsep=4pt, leftmargin=0em, itemindent=65pt, parsep=0pt]

\item\sloppy
Ввод функции <<Запрет АВР>> в работу осуществляется программным переключателем SGF1 <<Ввод\_функции>>. Информация о введенной функции отображается выходным сигналом <<ЛО НН / ЗАВР: Ввод>>.
\item
Функция <<Запрет АВР>> может быть выведена из работы оперативно путем активации сигнала <<ОВ ЛО НН / ЗАВР>>, что характеризуется активным состоянием сигнала <<ЛО НН / ЗАВР: Оперативный вывод>> на выходе функции.
\item
Функция принимает сигнал срабатывания от общей логики отключения и формирует сигналы <<ЛО НН / ЗАВР: Запрет АВР>>.
\item
Алгоритм работы <<Запрет АВР>> выполнен в соответствии с рисунком \ref{lonn:img_zavr}. В таблице \ref{lonn:tbl_azvr} приведены параметры, необходимые для настройки функции <<Запрет АВР>>.

\vspace{3mm}
\begin{figure}[H]
\centering

\includegraphics[width=0.9\textwidth,height=0.9\textheight,keepaspectratio]{\fbpath/61. ЛО В НН Т 35 Т2/_latex/img/zavr.pdf}
\captionof{figure}{Функциональная схема <<Запрет АВР>>}\label{lonn:img_zavr}
\end{figure}

\small
\begin{longtable}{|>{\centering\arraybackslash}m{5cm}|>{\centering\arraybackslash}m{4cm}|>{\centering\arraybackslash}m{4cm}|>{\centering\arraybackslash}m{2cm}|>{\centering\arraybackslash}m{1cm}|}
\caption{Параметры для настройки функции <<Запрет АВР>>\hfill\vspace{-0.5\baselineskip}}\label{lonn:tbl_azvr}\\ 
\hline
\rowcolor{gray!30}
Параметр (Параметр на ИЧМ) & Условное обозначение на схеме & Значение/ Диапазон & Единица измерения & Шаг \\ 
\hline
\endfirsthead
%\caption{\emph{Продолжение\hfill\vspace{-0.5\baselineskip}}} \\ % Отступ обозначения таблицы от таблицы и выравнивание надписи слева
\caption*{\hspace{3pt}\emph{Продолжение таблицы \ref{lonn:tbl_azvr}\hfill\vspace{-0.5\baselineskip}}} \\ % сделано по ГОСТ 2.105 п.6.8.7
\hline
\rowcolor{gray!30}
Параметр (Параметр на ИЧМ) & Условное обозначение на схеме & Значение/ Диапазон & Единица измерения & Шаг \\ 
%\hline
\endhead
%\hline
\endfoot
%\hline
\endlastfoot
%==+t1*LVBTSRBLC1|LVTRESCBOFF1
%===t1*LVBTSRBLC1
\centering Ввод функции в работу (Ввод\_функции) & \centering SGF1 & \centering 0 = Не предусмотрено\\1 = Предусмотрено & \centering -- & \centering \arraybackslash -- \\
\hline
%===t1
\end{longtable}
\normalsize

\end{enumerate}

\end{enumerate}